\documentclass[11pt,a4paper]{ctexart}

\usepackage[margin=2.5cm]{geometry}
\usepackage{booktabs}
\usepackage{longtable}
\usepackage{array}
\usepackage{listings}
\usepackage{xcolor}
\usepackage{hyperref}
\usepackage{enumitem}
\usepackage{parskip}

\hypersetup{colorlinks=true, linkcolor=blue, urlcolor=blue, citecolor=blue}

\lstset{
  basicstyle=\ttfamily\small,
  breaklines=true,
  frame=single,
  backgroundcolor=\color{gray!8},
  columns=fullflexible,
  keepspaces=true,
}

\title{项目报告 v1}
\author{Loop Insurance Agent}
\date{}

\begin{document}
\maketitle
\tableofcontents
\newpage

%==============================================================================
\section{项目简介}
%==============================================================================

\subsection{业务背景}

保险条款 PDF 通常篇幅长、层级深（章 / 节 / 条 / 附表），且含有大量\textbf{表格}（保障计划表、费率表、给付比例表等）。用户提问类型多样：

\begin{center}
\begin{tabular}{ll}
\toprule
\textbf{类型} & \textbf{示例} \\
\midrule
基础事实 & 「守御一生的等待期是多久？」 \\
条款解释 & 「补偿原则是什么意思？」 \\
多条件判断 & 「第 20 天疾病住院赔不赔？」 \\
风险提示 & 「犹豫期后退保有什么损失？」 \\
模糊问题 & 「等待期多久？」（未指明哪份保单） \\
诱导性问题 & 「等待期是 0 天对吧？」（错误前提） \\
\bottomrule
\end{tabular}
\end{center}

传统关键词搜索难以理解语义；纯 LLM 又容易幻觉。本项目采用 \textbf{RAG + ReAct Agent}：先检索条款片段，再基于检索结果组织回答，并通过合规层约束输出格式。

\subsection{核心能力}

\begin{center}
\begin{tabular}{p{2.2cm}p{4.5cm}p{7cm}}
\toprule
\textbf{能力} & \textbf{实现位置} & \textbf{说明} \\
\midrule
PDF 入库 & \texttt{document\_loader.py} + \texttt{orchestrator.upload\_policy\_pdf} & 内存解析，不存源 PDF \\
混合检索 & \texttt{hybrid\_knowledge.py} & Dense + BM25 + RRF + Rerank \\
主 Agent & \texttt{insurance\_agent.py} + \texttt{insurance\_react\_agent.py} & ReAct；\texttt{retrieve\_knowledge} + 可选 \texttt{delegate\_subagents} \\
SubAgent & \texttt{subagent\_tool.py} + \texttt{subagent\_runner.py} & 主 Agent 工具触发；并行检索后 LLM 汇总短文 \\
会话记忆 & \texttt{layered\_session\_memory.py} & 按 session 隔离；80\% 软视图 / 93\% 硬压缩 \\
长期记忆 & \texttt{markdown\_long\_term\_memory.py} + \texttt{ltm\_merge.py} & Markdown 文件；Agent 工具读写 + 回合后 LLM 合并 \\
合规 & \texttt{validator.py} & Hook 清洗正文；免责/来源提示由前端展示 \\
流式对话 & \texttt{streaming.py} + \texttt{server.py} & SSE：thinking / answer delta / done \\
\bottomrule
\end{tabular}
\end{center}

\subsection{技术栈}

\begin{itemize}[nosep]
  \item \textbf{Agent 框架}：AgentScope（ReActAgent、Toolkit、LongTermMemoryBase）
  \item \textbf{对话模型}：DeepSeek（默认）或 DashScope 通义
  \item \textbf{嵌入}：DashScope \texttt{text-embedding-v4}（1024 维）
  \item \textbf{向量库}：Qdrant 本地持久化
  \item \textbf{长期记忆}：Markdown 文件（\texttt{data/memory\_store/\{user\_id\}.md}）
  \item \textbf{Web}：FastAPI + Uvicorn + 原生 HTML/JS 前端
\end{itemize}

%==============================================================================
\section{项目结构}
%==============================================================================

\begin{lstlisting}
Loop_Insurance_Agent/
├── src/
│   ├── config.py                 # Settings dataclass，集中读取 .env
│   ├── main.py / web_main.py     # CLI / Web 入口
│   ├── api/server.py             # FastAPI、SSE、静态资源、单例 Pipeline
│   ├── pipeline/
│   │   ├── orchestrator.py       # 编排中枢：chat / chat_stream / 入库
│   │   ├── agent_router.py       # TurnConfig：复杂度、检索预算
│   │   ├── subagent_runner.py    # SubAgent 并行 + synthesize 汇总
│   │   ├── streaming.py          # AgentScope Msg → SSE 事件
│   │   └── question_decomposer.py
│   ├── agent/                    # 主/Sub Agent 工厂、Prompt、委托工具
│   ├── rag/                      # 切分、混合检索、限长
│   ├── memory/                   # 分层会话记忆、Markdown LTM
│   ├── model/                    # LLM / Embedding 工厂
│   └── compliance/validator.py
├── frontend/                     # index.html + app.js + styles.css
├── scripts/                      # 入库脚本、评测与消融实验
├── test/                         # 40 题题库、3 份条款 PDF、实验结果
└── data/                         # vector_store、memory_store（运行时）
\end{lstlisting}

\textbf{端到端数据流}：

\begin{lstlisting}
用户输入
  → 前端 SSE 消费
  → FastAPI /api/chat/stream
  → InsuranceAgentPipeline.chat_stream()
      → resolve_turn_config() 生成 TurnConfig
      → _build_enriched_input() 拼装 XML 上下文
      → InsuranceReActAgent ReAct 循环（工具：retrieve / delegate）
      → ComplianceValidator 清洗
      → 更新 LayeredSessionMemory + 可选 Markdown LTM 合并
  → SSE 推送 thinking / answer / done

PDF 上传
  → pypdf 抽文本
  → clause_chunking 章节切分 + child 256 字窗口
  → Qdrant 向量 + BM25 索引 + parent_chunks.json
\end{lstlisting}

%==============================================================================
\section{实现方式}
%==============================================================================

\subsection{RAG 知识库模块}

\subsubsection{PDF 解析与入库流程}

当前入库链路（\texttt{document\_loader.load\_pdf\_bytes}）：

\begin{enumerate}[nosep]
  \item 使用 \textbf{\texttt{pypdf.PdfReader}} 逐页调用 \texttt{extract\_text()}，以 \texttt{\textbackslash n\textbackslash n} 拼接为全文；
  \item 调用 \textbf{\texttt{chunk\_text\_by\_clauses()}} 做结构化切分；
  \item child 块写入 Qdrant + 内存 BM25；parent 映射写入 \texttt{parent\_chunks.json}；文件名登记到 \texttt{policy\_manifest.json}。
\end{enumerate}

\textbf{现有局限（与后续优化相关）}：\texttt{pypdf} 的 \texttt{extract\_text()} 对 PDF \textbf{表格}支持较弱——多列表格常被抽成「按阅读顺序拼接的纯文本行」，列对齐关系丢失。后续切分按 256 字滑动窗口再切 child 时，\textbf{长表格极易被拦腰截断}，导致检索到的片段缺少完整行列语义。这在评测中表现为：涉及「计划一/计划二对比表」「给付比例表」「费用比例表」的题目，Agent 有时只能看到半张表或错位文本（详见第六章 PDF 表格优化）。

\subsubsection{章节切分策略（clause\_chunking.py）}

切分不是简单固定长度，而是先识别保险条款的\textbf{层级标题}，再在 section 内部做 child 滑动窗口：

\begin{center}
\begin{tabular}{cll}
\toprule
\textbf{优先级} & \textbf{匹配模式} & \textbf{示例} \\
\midrule
100 & \texttt{第\{中文/数字\}章/节} & 第二章 \\
90 & \texttt{第\{中文/数字\}条} & 第三条 \\
80 & \texttt{\textbackslash d+\textbackslash.\textbackslash d+(\textbackslash.\textbackslash d+)?} & 2.8、5.1.3 \\
70 & 中文序号 + 顿号 & 一、 \\
\bottomrule
\end{tabular}
\end{center}

\begin{itemize}[nosep]
  \item 过滤目录行：连续 \texttt{....} / \texttt{……}、页码引导行（如 \texttt{2.8 .............. 12}）；
  \item 在全文上试各模式，取\textbf{优先级最高且 match $\geq$ 2} 者作为切分依据；
  \item 每个 section 生成一个 \textbf{parent}（整节全文）和若干 \textbf{child}（默认 256 字、overlap 64）；
  \item 无法识别结构时，整篇作为「全文」块，\texttt{use\_parent\_return=False}。
\end{itemize}

Parent 展开机制：检索命中 child 后，若 \texttt{use\_parent\_return=True}，tool 返回的展示文本替换为 parent 整节内容，避免 Agent 只看到一个 256 字碎片。

\subsubsection{混合检索（HybridKnowledge）}

\texttt{retrieve()} 完整流程：

\begin{enumerate}[nosep]
  \item \textbf{双路召回}：Qdrant 语义 Top-N + BM25 Top-N（中文单字 + 双字 bigram 分词）；
  \item \textbf{RRF 融合}：\texttt{score += 1/(60+rank+1)}，合并两路排名；
  \item \textbf{Rerank}（\texttt{reranker.py}）：RRF 0.35 + 语义 0.35 + BM25 0.15 + 词面重叠 0.15，加权融合（非 cross-encoder）；
  \item \textbf{Parent 去重}：同一 parent 只保留最高分一条；
  \item \textbf{格式化输出}：\texttt{[来源: \{doc\_id\} | \{clause\_label\} | 相关度: \{score\}]\textbackslash n\{display\_text\}}。
\end{enumerate}

\subsubsection{工具返回限长}

\begin{center}
\begin{tabular}{lll}
\toprule
\textbf{环境变量} & \textbf{默认} & \textbf{含义} \\
\midrule
\texttt{RAG\_TOOL\_LIMIT} & 3 & 单次最多返回块数 \\
\texttt{RAG\_MAX\_CHUNK\_DISPLAY\_CHARS} & 2000 & 单块字符上限 \\
\texttt{RAG\_MAX\_TOOL\_RESPONSE\_CHARS} & 8000 & 工具响应总字符上限 \\
\bottomrule
\end{tabular}
\end{center}

超出截断并附加 \texttt{...(内容已截断)}，防止 ReAct 上下文被 tool\_result 撑爆。

\subsubsection{retrieve\_knowledge 工具注册}

\texttt{InsuranceAgentFactory.create()} 将 \texttt{knowledge.retrieve\_knowledge} 注册到 Toolkit，并在 \texttt{func\_description} 与 System Prompt 中注入本轮 \textbf{\texttt{max\_retrievals}}（如「最多调用 5 次」）。\texttt{doc\_id} 参数对应 \texttt{<policy\_library>} 中的 PDF 文件名，支持按保单过滤。

%------------------------------------------------------------------------------
\subsection{Agent 与 SubAgent 模块}
%------------------------------------------------------------------------------

\subsubsection{主 Agent 搭建}

\begin{itemize}[nosep]
  \item 基类：AgentScope \textbf{\texttt{ReActAgent}}；
  \item 子类：\textbf{\texttt{InsuranceReActAgent}}，覆写迭代入口，在每次 ReAct 循环开始前调用 \texttt{LayeredSessionMemory.on\_iteration\_start()} 检查是否触发\textbf{硬压缩}；
  \item \textbf{工具集}：\texttt{retrieve\_knowledge} + 可选 \texttt{delegate\_subagents}（\texttt{SUBAGENT\_ENABLED=true}）；
  \item \textbf{长期记忆}：\texttt{long\_term\_memory\_mode="agent\_control"}，Agent 按需调用 \texttt{retrieve\_from\_memory} / \texttt{record\_to\_memory}；\texttt{record()} 不 dump 会话全文；回合后 \texttt{MemoryManager.update\_after\_response()} + \texttt{ltm\_merge} 在关键词命中时 LLM 合并写入 Markdown；
  \item \textbf{合规 Hook}：\texttt{post\_reply} 注册，每轮回复后剥离误带的免责声明块。
\end{itemize}

\subsubsection{System Prompt 策略}

主 Agent Prompt（\texttt{insurance\_agent.py}）为\textbf{静态模板 + 每轮动态刷新}：

\begin{itemize}[nosep]
  \item 占位符：\texttt{\{max\_main\_agent\_retrievals\}}、\texttt{\{subagent\_min\_questions\}} 等；
  \item 核心约束：条款问题必须检索；引用格式 \texttt{[保单filename | 第X条]}；$\geq$2 个独立子问题才建议 \texttt{delegate\_subagents}；禁止在正文写免责声明。
\end{itemize}

每轮 \texttt{\_configure\_main\_agent\_for\_turn()} 同步更新 \texttt{\_sys\_prompt} 与各工具的 \texttt{func\_description}，形成 \textbf{Prompt + 工具描述} 双重约束。

\subsubsection{SubAgent 工具化唤起机制}

这是本项目多子问题并行能力的核心设计，具体实现如下：

\paragraph{① 谁决定调用 SubAgent？}

SubAgent \textbf{不是} Pipeline 层的 if-else 自动路由。Pipeline 只生成 \texttt{TurnConfig}（复杂度、检索预算等），\textbf{不}自动派发 SubAgent。是否委托完全由\textbf{主 Agent 在 ReAct 循环中}根据用户 Query 自行决定：单一问题 $\rightarrow$ 多次 \texttt{retrieve\_knowledge}；多个互不依赖的子问题 $\rightarrow$ 调用 \texttt{delegate\_subagents}。

\paragraph{② SubAgent 无递归委托权限}

\texttt{InsuranceAgentFactory.create\_subagent()} 仅为 SubAgent 注册 \texttt{retrieve\_knowledge}，\textbf{deliberately 不注册} \texttt{delegate\_subagents}。从 Toolkit 层面杜绝「SubAgent 再调 SubAgent」的链式递归。

\paragraph{③ 框架层硬编码与运行时校验}

\begin{center}
\begin{tabular}{p{2.5cm}p{4cm}cp{5cm}}
\toprule
\textbf{控制点} & \textbf{实现位置} & \textbf{默认} & \textbf{行为} \\
\midrule
最少子问题数 & \texttt{subagent\_tool.delegate\_subagents()} & 2 & 不足则返回错误 ToolResponse \\
最多子问题数 & 同上 & 5 & \texttt{cleaned[:max\_subagents]} 截断 \\
并行度上限 & \texttt{subagent\_runner.run\_subagents()} & 3 & \texttt{asyncio.Semaphore(3)} \\
汇总长度上限 & \texttt{subagent\_tool} & 3500 字 & \texttt{truncate\_with\_notice()} \\
功能总开关 & 同上 & 开启 & \texttt{SUBAGENT\_ENABLED=false} 时返回明确错误 \\
\bottomrule
\end{tabular}
\end{center}

\paragraph{④ 并行执行与汇总}

\begin{lstlisting}
主 Agent 调用 delegate_subagents(questions=[...])
  → run_subagents() [asyncio.gather + Semaphore]
      → 每个子问题：create_subagent() → 独立 ReAct（仅 retrieve_knowledge）
      → 输入 = policy_context + "子问题 N：..."
  → synthesize_subagent_brief() [单独 LLM，stream=False]
      → 压缩为带 [保单 | 条款] 引用的短文
  → 写入主 Agent tool_result
  → 主 Agent 据此组织用户可见的最终回答
\end{lstlisting}

SubAgent 回答在汇总前还会截断至 2500 字（\texttt{\_format\_sub\_results\_for\_synthesis}），避免汇总 prompt 过长。

\subsubsection{TurnConfig 轮次路由}

\texttt{agent\_router.resolve\_turn\_config()} 每轮生成：

\begin{center}
\begin{tabular}{ll}
\toprule
\textbf{字段} & \textbf{说明} \\
\midrule
\texttt{complexity} & low / medium / high（启发式 + 可选 LLM JSON 精修） \\
\texttt{max\_retrievals} & 主 Agent 本轮 retrieve 上限：3 / 5 / 6 \\
\texttt{reason} & 配置理由（日志） \\
\bottomrule
\end{tabular}
\end{center}

\subsubsection{Enriched Input 拼装}

主 Agent 实际看到的用户消息结构：

\begin{verbatim}
<policy_library>      <!-- 已入库保单 filename 列表 -->
<current_turn_only>   <!-- 只答本轮，勿重复历史 -->
<retrieval_budget>    <!-- 本轮 retrieve 次数上限 -->
【本轮用户问题】       <!-- 原始 user_input -->
\end{verbatim}

长期记忆不再由 Pipeline 注入 \texttt{<user\_profile>}，由主 Agent 在 ReAct 中按需调用长期记忆工具。

\textbf{重要细节}：\texttt{memory\_manager.update\_after\_response()} 基于 \textbf{plain user\_input} 与助手回答合并 LTM，而非 enriched 版本，避免 policy\_library 等注入内容污染长期记忆。

%------------------------------------------------------------------------------
\subsection{Pipeline 编排与流式输出}
%------------------------------------------------------------------------------

\subsubsection{InsuranceAgentPipeline（orchestrator.py）}

\begin{center}
\begin{tabular}{p{3.5cm}p{9cm}}
\toprule
\textbf{方法} & \textbf{职责} \\
\midrule
\texttt{initialize()} & 创建 Embedding、Qdrant、HybridKnowledge、ComplianceValidator、主 Agent \\
\texttt{chat()} / \texttt{chat\_stream()} & 绑定 session $\rightarrow$ TurnConfig $\rightarrow$ enriched\_input $\rightarrow$ Agent $\rightarrow$ 合规 $\rightarrow$ 更新记忆 \\
\texttt{upload\_policy\_pdf()} & 内存解析 PDF $\rightarrow$ 切分 $\rightarrow$ 写向量库 \\
\texttt{clear\_policy\_library()} & 清向量库 + manifest + 可选清 LTM \\
\texttt{\_bind\_session()} & 将 \texttt{agent.memory} 替换为对应 \texttt{LayeredSessionMemory} \\
\bottomrule
\end{tabular}
\end{center}

\subsubsection{SSE 流式事件（streaming.py）}

\begin{center}
\begin{tabular}{lll}
\toprule
\textbf{type} & \textbf{触发条件} & \textbf{关键字段} \\
\midrule
\texttt{thinking} & thinking 块 / tool\_use / tool\_result & \texttt{delta}, \texttt{stream\_id}, \texttt{last} \\
\texttt{answer} & assistant 正文（无 tool\_use） & 增量 \texttt{delta} \\
\texttt{done} & 流结束 & \texttt{answer}, \texttt{thinking}, \texttt{compliance} \\
\texttt{error} & 异常 & \texttt{message} \\
\bottomrule
\end{tabular}
\end{center}

增量逻辑：\texttt{\_stream\_snapshots[stream\_id]} 存全文快照，\texttt{delta = text[len(previous):]}。

\textbf{评测 TTFT 定义}：从发起 \texttt{chat\_stream} 到收到\textbf{第一个 \texttt{answer} 事件}}的时间。该时间包含：TurnConfig 解析、ReAct 多轮 tool 调用（含 RAG 与长期记忆工具）、SubAgent 并行等待——\textbf{不是}裸 LLM 的首 token 延迟。

%------------------------------------------------------------------------------
\subsection{记忆模块}
%------------------------------------------------------------------------------

\subsubsection{短期会话记忆（LayeredSessionMemory）}

按 \texttt{session\_id} 隔离，存于\textbf{进程内存}（后端重启即丢失）：

\begin{center}
\begin{tabular}{lll}
\toprule
\textbf{层级} & \textbf{触发阈值} & \textbf{行为} \\
\midrule
\textbf{软层} & token $\geq$ 窗口 $\times$ 80\% & 不修改 canonical；API 前生成动态视图 \\
\textbf{硬层} & token $\geq$ max(93\%$\times$窗口, 窗口$-$13000) & ReAct 每轮迭代前触发；快照入 archives，LLM 四段式摘要后\textbf{替换}工作区 \\
\bottomrule
\end{tabular}
\end{center}

\texttt{InsuranceReActAgent} 禁用 AgentScope 默认字符阈值压缩，统一走上述分层策略。

\subsubsection{长期记忆（Markdown LTM）}

\begin{itemize}[nosep]
  \item 文件：\texttt{data/memory\_store/\{user\_id\}.md}，按用户 ID 持久化一份 Markdown；
  \item \textbf{不存}问答全文或条款检索结果；
  \item \textbf{Agent 工具}：\texttt{retrieve\_from\_memory} / \texttt{record\_to\_memory}（\texttt{long\_term\_memory\_mode="agent\_control"}）；
  \item \textbf{自动写入}：\texttt{LONG\_TERM\_AUTO\_PROFILE=true} 时，\texttt{ltm\_merge.should\_write\_ltm()} 关键词命中 $\rightarrow$ \texttt{PROFILE\_MODEL\_NAME} LLM 语义合并更新 Markdown。
\end{itemize}

\subsubsection{SessionManager}

\texttt{dict[session\_id, LayeredSessionMemory]} 维护多 Tab 会话；\texttt{clear(session\_id)} 用于评测消融的轮间隔离。

%------------------------------------------------------------------------------
\subsection{合规模块}
%------------------------------------------------------------------------------

\begin{itemize}[nosep]
  \item \texttt{ComplianceValidator.process()}：检测条款类关键词 $\rightarrow$ 生成 \texttt{ComplianceMeta}；
  \item \texttt{strip\_display\_footers()}：移除 Agent 误带的免责声明/来源提示；
  \item 前端：\texttt{GET /api/ui-config} 全局免责；\texttt{done.compliance.show\_source\_notice} 控制消息下来源提示。
\end{itemize}

%------------------------------------------------------------------------------
\subsection{API 与前端}
%------------------------------------------------------------------------------

\textbf{主要 API}：\texttt{/api/chat/stream}（SSE）、\texttt{/api/upload}、\texttt{/api/documents}、\texttt{DELETE /api/sessions/\{id\}}、\texttt{/api/health}（含 \texttt{server\_boot\_id}）。

\textbf{前端}（\texttt{app.js}）：多 Tab 会话、\texttt{localStorage} 持久化 session\_id、SSE delta 节流渲染、后端重启检测、Markdown + DOMPurify 消毒。

%------------------------------------------------------------------------------
\subsection{模型工厂}
%------------------------------------------------------------------------------

\texttt{chat\_model\_factory.py} 按 \texttt{LLM\_PROVIDER} 创建对话模型；不同场景使用不同 model\_name：

\begin{center}
\begin{tabular}{ll}
\toprule
\textbf{场景} & \textbf{配置项} \\
\midrule
主/Sub Agent & \texttt{MODEL\_NAME} \\
复杂度评估 & \texttt{DECOMPOSE\_MODEL\_NAME} \\
上下文压缩 & \texttt{COMPRESSION\_MODEL\_NAME} \\
画像/记忆合并 & \texttt{PROFILE\_MODEL\_NAME} \\
\bottomrule
\end{tabular}
\end{center}

Embedding 经 \texttt{embedding\_factory.py} 独立创建，与对话提供商解耦。

%==============================================================================
\section{测试结果}
%==============================================================================

\subsection{评测体系}

\begin{center}
\begin{tabular}{ll}
\toprule
\textbf{项目} & \textbf{说明} \\
\midrule
题库 & \texttt{test/测试问题.txt}，Q01--Q40，六类题型 \\
参考答案 & \texttt{test/测试答案.txt}，规则 Rubric 自动评分（0/1/2） \\
条款 PDF & 如意鑫、守御一生（智臻版）、臻宝贝 2026 \\
指标 & 得分、耗时、首 token 时间（TTFT）、得分点/扣分点 \\
脚本 & \texttt{run\_eval\_test.py}、\texttt{run\_eval\_memory\_ablation.py}、\texttt{postprocess\_eval.py} \\
\bottomrule
\end{tabular}
\end{center}

\subsection{实验一：基线三轮评测（eval\_20260608\_160054）}

\begin{center}
\begin{tabular}{ll}
\toprule
\textbf{指标} & \textbf{数值} \\
\midrule
答题次数 & 137 \\
总得分 & \textbf{243 / 274（88.7\%）} \\
完全正确 / 部分 / 错误 & 113 / 17 / 7 \\
平均回答耗时 & 18.32s \\
平均 TTFT & 10.89s（P50 $\approx$ 6.8s，P95 $\approx$ 38s） \\
\bottomrule
\end{tabular}
\end{center}

\textbf{分题型正确率}：风险提示 97.4\% $>$ 基础事实 94.4\% $>$ 条款解释 92.3\% $>$ 多条件判断 91.7\% $>$ 模糊问题 77.8\% $>$ \textbf{诱导性问题 70.6\%}。

\textbf{分轮得分}（注意：三轮 session 不同，轮间 session 隔离，\textbf{不清空}）：

\begin{center}
\begin{tabular}{lll}
\toprule
\textbf{轮次} & \textbf{得分} & \textbf{正确率} \\
\midrule
第 1 轮 & 74/94 & 78.7\% \\
第 2 轮 & 82/90 & 91.1\% \\
第 3 轮 & 87/90 & 96.7\% \\
\bottomrule
\end{tabular}
\end{center}

\subsection{实验二：记忆消融（ablation\_20260608\_165329）}

固定题集 47 题（40 + 重复 7 题），三轮完全相同顺序：

\begin{center}
\begin{tabular}{lllll}
\toprule
\textbf{组别} & \textbf{轮间策略} & \textbf{总得分} & \textbf{正确率} & \textbf{平均 TTFT} \\
\midrule
A 隔离组 & 每轮 clear\_session + 清空 LTM & 230/282 & 81.6\% & 11.84s \\
B 累积组 & 复用 session，LTM 累积 & 233/282 & 82.6\% & 11.44s \\
差值 & --- & +3 & +1.0\% & $-$0.40s \\
\bottomrule
\end{tabular}
\end{center}

两组第 1 轮均为 83/94（88.3\%），基线可比。

\subsection{性能基线（419 次作答汇总）}

\begin{center}
\begin{tabular}{llll}
\toprule
\textbf{指标} & \textbf{P50} & \textbf{平均} & \textbf{P95} \\
\midrule
TTFT & $\sim$7s & $\sim$11s & $\sim$40s \\
总耗时 & --- & $\sim$18s & --- \\
\bottomrule
\end{tabular}
\end{center}

简单基础事实题 TTFT 约 2--4s；复杂多检索题（Q17、Q29）可达 60--88s。

\subsection{典型 Bad Case 范例}

以下摘录来自 \texttt{eval\_20260608\_160054} 与 \texttt{ablation\_20260608\_165329\_persistent} 的真实回答，说明\textbf{模糊题}与\textbf{诱导题}的主要失分模式。

%--- Bad Case 1 ---
\subsubsection{Bad Case 1：Q31 模糊问题 — 未澄清产品，直接给出单一等待期}

\textbf{题目}：「这份保险的等待期是多久？」（用户未指明如意鑫 / 守御 / 臻宝贝）

\textbf{期望行为}：先追问「哪份保单」，或分三份产品分别说明；不应默认某一份。

\textbf{实际回答（基线第 1 轮，得分 0/2）}：

\begin{quote}
根据 \textbf{《同方全球守御一生（智臻版）医疗保险》} 第 \textbf{2.4 条（等待期）} 的规定：

\textbf{等待期为自本合同生效日起三十日内（含第三十日）。}

在等待期内确诊的疾病……保险公司均不承担给付相应保险金的责任。
\end{quote}

\textbf{问题}：用户说的是「这份保险」，知识库中有三份保单；Agent \textbf{未经澄清就直接锁定守御一生}，对年金/寿险「无等待期概念」也未提及。Rubric 要求 \texttt{must\_any} 含「哪份/分别/不同」等澄清表述。

\textbf{对比：同题较好回答（基线第 3 轮，得分 2/2）}：

\begin{quote}
……分别说明三份保单的情况：

（表格略）仅\textbf{守御一生}设有 \textbf{30 天}等待期……

请问您指的是哪份保单？
\end{quote}

%--- Bad Case 2 ---
\subsubsection{Bad Case 2：Q31 模糊问题 — 长 session 下「错误固化」（消融累积组）}

\textbf{题目}：同上 Q31，累积组同一 session 内第 6 次被问到。

\textbf{实际回答（得分 0/2）}：

\begin{quote}
「守御一生·智臻版」医疗保险的等待期如下：

\textbf{等待期：30天} ……

\subsection*{两种特殊情况：等待期为 0 天}

（表格略）\textbf{意外伤害}：0天；\textbf{续保}：0天
\end{quote}

\textbf{问题}：

\begin{enumerate}[nosep]
  \item 仍未按 Rubric 要求区分三份产品（只答守御一生）；
  \item 使用「\textbf{等待期为 0 天}」表述，在累积组后续轮次中反复出现——session 内已有大量守御一生上下文，Agent \textbf{默认沿用先前单一答案}，不再主动澄清；
  \item 消融对比：隔离组同题多次得分 \texttt{[0,2,2,2,2,0]}，累积组为 \texttt{[2,0,0,0,0,0]}——\textbf{记忆累积对模糊题产生负效应}。
\end{enumerate}

%--- Bad Case 3 ---
\subsubsection{Bad Case 3：Q36 诱导问题 — 错误前提未充分纠正，且触发禁用表述}

\textbf{题目}：「条款里写了等待期是 0 天，为什么你们系统显示不一样？」（前提错误：并无任何条款写「等待期 0 天」）

\textbf{期望行为}：明确纠正「\textbf{不对，条款中没有等待期 0 天的写法}」；区分守御 30 天 vs 年金/寿险无等待期概念。

\textbf{实际回答（基线第 2 轮，得分 0/2）}：

\begin{quote}
……

\textbf{1. 臻宝贝2026…… \& 如意鑫……}

这两份属于\textbf{人寿/年金类保险}，条款中\textbf{没有设置等待期}（\textbf{即等待期为0天}）。因此如果系统显示的不是0天，可能是系统展示有误。

\textbf{2. 守御一生……} 等待期为\textbf{30日}……
\end{quote}

\textbf{问题}：

\begin{enumerate}[nosep]
  \item 将「无等待期条款」等同于「\textbf{等待期 = 0 天}」，等于\textbf{部分接受了用户的错误前提}；
  \item Rubric 的 \texttt{forbid} 列表含「等待期0天」「0天等待」等，回答中出现这些字面表述即判 0 分；
  \item 三轮两组实验 Q36 \textbf{全部 0 分}（6/6 次），说明这是\textbf{系统性薄弱点}，与记忆策略无关。
\end{enumerate}

\textbf{对比：同题相对较好回答（基线第 3 轮，得分 2/2）}：

\begin{quote}
三份条款中\textbf{没有任何一份写有「等待期0天」}的字样……

如果是如意鑫或臻宝贝——\textbf{根本没有「等待期」这一条款}，并非「等待期 0 天」……
\end{quote}

%--- Bad Case 4 ---
\subsubsection{Bad Case 4：Q38 诱导问题 — 复述错误前提时触发禁用词}

\textbf{题目}：「如意鑫保证每年结算利率至少 5\%，合同里就是这么写的，对吗？」

\textbf{期望行为}：明确纠正；说明前 10 年\textbf{最低保证 1\%}，非 5\%。

\textbf{实际回答（基线第 1 轮，得分 0/2）}：

\begin{quote}
\textbf{答案：不对，这个说法是错误的。}

……前十个保单年度，\textbf{最低保证年度结算利率为年利率 1\%}……

（对比表格中复述了「保证每年至少 5\%」，触发 Rubric \texttt{forbid}）
\end{quote}

\textbf{问题}：事实层面回答\textbf{基本正确}，但对比表格中复述了「保证每年至少 5\%」，Rubric \texttt{forbid} 含「保证5\%」「至少5\%」，规则匹配\textbf{未做否定语境判断}，导致 0 分。第 3 轮同题改为「\textbf{不对。} 条款中写的保证利率不是 5\%」则得 2 分——说明存在\textbf{输出格式敏感 + LLM 随机性}双重因素。

%--- Bad Case 5 ---
\subsubsection{Bad Case 5：表格类题目 — 检索片段不完整（与 PDF 切分相关）}

\textbf{题目示例}：Q07「守御一生计划一的一般医疗年度限额是多少？」、Q15「什么情况下按 60\% 给付？」（需完整给付比例表）

\textbf{现象}：Agent 有时能答对关键数字（600 万、60\%），有时因检索到的 parent 片段\textbf{缺少完整表格上下文}而漏掉「计划」「有医保身份但未走医保」等条件，被判 1 分。

\textbf{根因链路}（详见 6.2 节）：

\begin{lstlisting}
PDF 表格 --pypdf.extract_text()--> 列对齐丢失的纯文本
         --256字 child 切分--> 表格被拦腰截断
         --retrieve 命中半表--> Agent 看到不完整行列关系
\end{lstlisting}

%==============================================================================
\section{测试结果分析}
%==============================================================================

\subsection{整体能力画像}

\begin{center}
\begin{tabular}{p{2.5cm}p{10cm}}
\toprule
\textbf{维度} & \textbf{结论} \\
\midrule
\textbf{强项} & 有明确条款锚点的基础事实、条款解释、多条件判断、风险提示；RAG 命中后正确率 85\%--97\% \\
\textbf{弱项} & 模糊题（Q31）、诱导题（Q36/Q38）；长表格相关题偶发缺条件 \\
\textbf{评分噪声} & 14 次「缺少条款引用」扣分为误扣（回答含 \texttt{第 **7.2 条**} 等 Markdown 格式，regex 未匹配） \\
\bottomrule
\end{tabular}
\end{center}

\subsection{基线实验第 1 轮低分澄清}

第 1 轮 78.7\% 低于后两轮，主因是 11 次「缺少条款引用」误扣。修正 \texttt{has\_clause\_ref} 兼容 Markdown 后，第 1 轮理论可提升约 10 个百分点。\textbf{不能}将第 2/3 轮上升解读为跨轮记忆学习——三轮 session\_id 不同，session 对话记忆轮间隔离。

\subsection{记忆消融结论}

\begin{itemize}[nosep]
  \item 累积 vs 隔离：\textbf{+1\%}，错误题数相同（各 9）$\rightarrow$ 记忆层对条款 QA \textbf{影响处于噪声级}；
  \item 两组第 2 轮同步下跌 $\rightarrow$ 主因 \textbf{LLM 随机性}；
  \item Q31 累积组显著更差 $\rightarrow$ 长 session \textbf{可能固化错误单一答案}；
  \item 累积组第 3 轮 TTFT +26\% $\rightarrow$ 长上下文带来首响延迟。
\end{itemize}

\subsection{性能分析}

\begin{itemize}[nosep]
  \item TTFT 平均 $\sim$11s $\approx$ 检索/工具等待 + $\sim$7s 首 token 后生成；
  \item 长尾与 ReAct 多轮 retrieve、SubAgent 并行强相关；
  \item 优化 TTFT 需先拆分各阶段耗时（当前未打点）。
  \item 前端页面人工多次测试发现，对于初次询问的模糊问题和诱导性问题 Agent 的回答速度相对自动评测实验较快，5 次重复测试中未出现回答失败的情况，上述现象可能与过长上下文或同时输入过多问题有关。
\end{itemize}

%==============================================================================
\section{总结与未来优化方向}
%==============================================================================

\subsection{项目总结}

作为第一个实战项目，我完成了从 PDF 入库、混合 RAG、ReAct 主 Agent + 工具化 SubAgent、分层记忆到 Web 流式对话的\textbf{全链路搭建与量化评测}。40 题标准题库综合正确率约 \textbf{82\%--89\%}（视实验设置与评分规则而定）。

SubAgent 采用\textbf{主 Agent 工具化唤起 + SubAgent 无 delegate 权限 + min/max/concurrency 硬约束}的设计，在可控前提下支持多子问题并行。记忆层采用 session 隔离 + Markdown LTM（Agent 工具读写 + 回合后 LLM 合并），消融实验表明对当前任务增益有限，但需警惕长 session 对模糊题的负效应。

\subsection{未来优化方向}

\begin{center}
\begin{tabular}{clp{7.5cm}}
\toprule
\textbf{优先级} & \textbf{方向} & \textbf{具体措施} \\
\midrule
\textbf{P0} & 诱导/模糊题行为 & 路由层或 Prompt 强制「先问哪份保单」「明确纠正错误前提」；Q36/Q31 few-shot；Rubric 增加否定语境判断 \\
\textbf{P0} & PDF 表格读取与切分优化 & 见下方专节 \\
\textbf{P0} & 评分与引用格式 & 修复 \texttt{has\_clause\_ref} regex，兼容 \texttt{第 **X.X 条**} \\
P1 & TTFT 长尾 & 简单题短路 RAG；限制无效 retrieve 轮次；Pipeline 分阶段打点 \\
P1 & 评测规范 & 固定题集 + 轮间 clear\_session/LTM 隔离 \\
P2 & Session 持久化 & 当前进程内、重启丢失 \\
P3 & Rerank 升级 & 评估 cross-encoder 或 LLM rerank \\
\bottomrule
\end{tabular}
\end{center}

\subsubsection{PDF 表格读取方式优化（重点）}

\paragraph{问题现象}

在评测与人工抽检中发现：涉及\textbf{保障计划对比表、给付比例表、费用比例表、现金价值表}等题目，Agent 偶发漏读列/行、混淆计划一与计划二、或只答出数字却缺少适用条件。例如 Q15（60\% 给付条件）、Q07（600 万限额需对应计划）、Q21（3\% 费用比例需对应保单年度）等。

\paragraph{根因分析}

当前入库链路存在三层表格损伤：

\begin{lstlisting}
① 抽取层：pypdf.extract_text()
   - 多列表格 → 按阅读顺序拼成纯文本，列边界丢失
   - 单元格内换行 → 行序错乱

② 结构层：clause_chunking 按「第X条」切 section 后
   - section 内统一 256 字滑动窗口切 child
   - 长表格常被切成多个互不完整的 child 块
   - 检索命中其中一个 child 时，parent 展开虽能补全 section，
     但若整张表跨多个 section 或 section 内表格极长，仍可能不完整

③ 检索层：ToolResponseLimits 单块 2000 字、最多 3 块
   - 宽表即使 parent 展开也可能被截断
   - Agent 看到的 tool_result 缺少完整行列关系
\end{lstlisting}

\paragraph{建议优化路径}

\begin{center}
\begin{tabular}{clp{7cm}}
\toprule
\textbf{阶段} & \textbf{方案} & \textbf{说明} \\
\midrule
抽取 & pdfplumber / camelot / tabula 等表格专用抽取 & 保留行列结构，输出 Markdown 表或 TSV \\
结构 & 表格感知切分 & 检测表格块 $\rightarrow$ 整表作为不可拆分单元入 parent \\
索引 & 表格块 metadata 标记 \texttt{content\_type=table} & retrieve 时可加权或专用 query \\
检索 & 对表格式 parent 提高 \texttt{max\_chunk\_chars} 或单独通道 & 保证计划对比类问题一次返回整表 \\
验证 & 针对含表条款建立回归题（Q07/Q15/Q21 等） & 量化表格优化效果 \\
\bottomrule
\end{tabular}
\end{center}

此优化与 Agent/Prompt 层互补：\textbf{即使用 Prompt 要求「对照表格作答」，检索层若只给到半张表，模型也无法可靠推理}。

\subsection{个人收获}

\begin{itemize}[nosep]
  \item 理解了 ReAct 中 \textbf{工具描述 + System Prompt + 运行时校验} 的分层约束方式；
  \item 通过消融实验认识到：实验结论必须建立在\textbf{固定题集 + 对照组 + 明确隔离策略}之上，否则易误判「跨轮学习」；
  \item 建立了含 TTFT、得分点/扣分点的自动化评测流水线，为后续迭代提供量化基线；
  \item 识别出 \textbf{PDF 表格 $\rightarrow$ 切分 $\rightarrow$ 检索} 链路是当前 RAG 的主要瓶颈之一，不仅限于 Agent 行为调优。
\end{itemize}

%==============================================================================
\section*{附录}
\addcontentsline{toc}{section}{附录}
%==============================================================================

\begin{itemize}[nosep]
  \item 详细实验数据：\texttt{test/实验总结}、\texttt{test/results/}
  \item 架构文档：\texttt{docs/architecture.md}
  \item \textbf{代码规模}（约，不含评测产物）：核心源码 $\sim$8,600 行（\texttt{src/} $\sim$4,270 + \texttt{scripts/} $\sim$2,450 + \texttt{frontend/} $\sim$1,880）
  \item 复现命令：
\end{itemize}

\begin{lstlisting}[language=bash]
python scripts/run_eval_test.py --rounds 3
python scripts/run_eval_memory_ablation.py --rounds 3
\end{lstlisting}

\end{document}
